% Однофазная нулевая (двухполупериодная) схема на тиристорах с VD0
\documentclass[border=4pt]{standalone}
\usepackage{amsmath}
\usepackage{fontspec}
\setmainfont{Calibri}
\usepackage[european, RPvoltages]{circuitikz}
\definecolor{ink}{RGB}{31,42,51}
\definecolor{grp}{RGB}{46,110,140}
\ctikzset{bipoles/length=1.0cm}
\ctikzset{bipoles/thickness=0.9}
\ctikzset{diodes/fill=ink}
\ctikzset{inductor=american}

\begin{document}
\begin{tikzpicture}[line width=1.0pt, color=ink, font=\small]
  % --- вторичная обмотка со средней точкой: w1 (верхняя), w2 (нижняя)
  \draw (0,3.4) to[L] (0,1.4);
  \draw (1.8,1.4) to[L] (1.8,-0.6);
  % --- средняя точка N и шина минуса (угол обмотки a — без точки)
  \draw (0,1.4) -- (1.8,1.4) node[circ]{} -- (8.2,1.4);
  \node[anchor=south west] at (0.55,1.44) {$N$};
  % --- тиристор VS1 (от верхней полуобмотки к шине +, один путь)
  \draw (0,3.4) to[Ty, name=VS1] (0,5.05) node[circ]{};
  % --- тиристор VS2 (от нижней полуобмотки, обход слева)
  \draw (1.8,-0.6) -- (1.8,-1.35) -- (-1.6,-1.35)
        -- (-1.6,3.95) to[Ty, name=VS2] (-1.6,4.75) -- (-1.6,5.05) node[circ]{};
  % --- шина плюса
  \draw (-1.6,5.05) -- (8.2,5.05);
  % --- обратный (нулевой) диод VD0: катод к шине +
  \draw (6.9,5.05) node[circ]{} -- (6.9,3.75);
  \draw (6.9,2.95) to[D] (6.9,3.75);
  \draw (6.9,2.95) -- (6.9,1.4) node[circ]{};
  % --- нагрузка
  \draw (8.2,5.05) -- (8.2,3.7) to[R] (8.2,1.75) -- (8.2,1.4);
  % --- подписи
  \node[anchor=south west] at (0.08,3.42) {$a$};
  \node[anchor=north west] at (1.88,-0.62) {$b$};
  \node[anchor=west] at (0.3,4.35) {VS1};
  \node[anchor=west] at (-1.25,4.35) {VS2};
  \node[anchor=east] at (6.5,3.35) {VD0};
  \node[anchor=east] at (7.75,3.2) {$u_d$};
  \node[anchor=west] at (8.65,2.7) {$R_{\text{н}}$};
  \node[anchor=south west] at (8.15,5.07) {$+$};
  \node[anchor=north west] at (8.15,1.38) {$-$};
  % --- вентили управления (подпись у затвора VS2)
  \node[anchor=north east, text=grp, font=\footnotesize] at (VS2.G) {от СИФУ};
\end{tikzpicture}
\end{document}
